% Two replacements in Online Appendix A.10. (OA.50)-(OA.52), the regularity
% conditions, and (OA.54)-(OA.56) stand. No \NUM placeholders.

% REPLACES: online_appendix.tex lines 1264-1271 ("When the reserve is above the entire low band ..." to "... because \Delta_T\mu>0.")
When the reserve is above the entire low band but below \(r\),
\(t_L=t_0\), \(g_L=0\), and the high-class formulas are unchanged.
Conditional on class and flow, exact \(V\) retains its within-class
distribution, so the benchmark price and residual proof applies to these
class-averaged payoffs whenever \(\Delta_T>0\). Entry is at least
\(\rho\) at every posterior: the zero-cost realization prepares because
class gross profit is nonnegative at each belief in this region and
positive at each belief above zero. In this region \(t_L=t_0\) is permitted: the price
construction remains strictly increasing on the bounded posterior
interval because \(\Delta_T\mu>0\).

% REPLACES: online_appendix.tex lines 1281-1299 ("the conditional order distributions, ..." to "... entry verification at such reserves.")
the conditional order distributions, the price mapping on flow, the
buyer's beliefs at every realized price including price atoms, and the
preparation rule at every price and cost. Its beliefs are the public
price posterior, not necessarily the raw order-flow posterior. Under the
declared two-point cost law entry is at least \(\rho\) at every flow, so
the inversion denominator cannot vanish through missing entry. A price
atom can still arise where the raw posterior is constant, as on a Laplace
plateau. The buyer then conditions on the whole preimage of that price,
whose members share the one plateau posterior. The construction used in C.6,
which prices each flow by the benchmark inversion wherever
\(\Delta_T>0\), is therefore the only pricing rule at its orders on that
part of the domain, up to those constant-posterior atoms.
\hyperref[oa-a-pools]{Section A.11} states this point and the reason a wider
family of pooling regions at fixed orders cannot arise under the zero cost.
Two continuations with the same orders but different price mappings would
still be different elements of
\(\mathcal E(p,r)\), and the numerical record identifies a continuation
by the whole object \eqref{eq:oa-oa-53}. The whole-object rule remains necessary,
because distinct order profiles occur. If \(\Delta_T=0\) or no sale is possible, the informative-price argument cannot be imported.
\hyperref[oa-c-reserve]{Section C.6} therefore requires direct pricing,
posterior, and entry verification at such reserves.
